\documentclass[11pt,a4paper]{article}
\usepackage[margin=2.2cm]{geometry}
\usepackage{amsmath,amssymb,mathtools}
\usepackage[T1]{fontenc}
\usepackage[utf8]{inputenc}
\usepackage{lmodern}
\usepackage{hyperref}

\title{Costanti Quartiche Watson\\Run Multi-modo ($n_{\mathrm{modes}}=2$), Ordine 2}
\author{}
\date{\today}

\begin{document}
\maketitle

\section*{Contesto}
Formule ottenute dal run:
\begin{center}
\texttt{python3 derive\_watson\_quartic\_vanvleck.py --max-order 2 --n-modes 2}
\end{center}
con notazione interna del codice per i coefficienti \texttt{mu1\_*} e frequenze \(\omega_0,\omega_1\).

\section*{Costanti principali}
\begin{align}
D_J &= \frac{1}{16}\left(\frac{X_0}{\omega_0^2}+\frac{X_1}{\omega_1^2}\right),\\[4pt]
D_{JK} &= \frac{1}{16}\left(\frac{Y_0}{\omega_0^2}+\frac{Y_1}{\omega_1^2}\right),\\[4pt]
D_K &= \frac{1}{16}\left(\frac{Z_0}{\omega_0^2}+\frac{Z_1}{\omega_1^2}\right),\\[4pt]
d_1 &= \frac{1}{16}\left(\frac{W_0}{\omega_0^2}+\frac{W_1}{\omega_1^2}\right),\\[4pt]
d_2 &= \frac{1}{16}\left(\frac{V_0}{\omega_0^2}+\frac{V_1}{\omega_1^2}\right).
\end{align}

\section*{Definizioni dei blocchi}
\begin{align}
X_0 &= 2\mu1_0\mu1_8+\mu1_2^2+2\mu1_2\mu1_6+\mu1_6^2,\\
X_1 &= 2\mu1_1\mu1_9+\mu1_3^2+2\mu1_3\mu1_7+\mu1_7^2.
\end{align}

\begin{align}
Y_0 &= 2\mu1_0\mu1_{16}-4\mu1_0\mu1_8+\mu1_{10}^2+2\mu1_{10}\mu1_{14}+\mu1_{12}^2+2\mu1_{12}\mu1_4+\mu1_{14}^2 \\
&\quad +2\mu1_{16}\mu1_8-2\mu1_2^2-4\mu1_2\mu1_6+\mu1_4^2-2\mu1_6^2,\\
Y_1 &= 2\mu1_1\mu1_{17}-4\mu1_1\mu1_9+\mu1_{11}^2+2\mu1_{11}\mu1_{15}+\mu1_{13}^2+2\mu1_{13}\mu1_5+\mu1_{15}^2 \\
&\quad +2\mu1_{17}\mu1_9-2\mu1_3^2-4\mu1_3\mu1_7+\mu1_5^2-2\mu1_7^2.
\end{align}

\begin{align}
Z_0 &= 2\mu1_{16}^2 - X_0 + Y_0,\\
Z_1 &= 2\mu1_{17}^2 - X_1 + Y_1.
\end{align}

\begin{align}
W_0 &= 2\mu1_0\mu1_{16}-\mu1_{10}^2-2\mu1_{10}\mu1_{14}+\mu1_{12}^2+2\mu1_{12}\mu1_4-\mu1_{14}^2-2\mu1_{16}\mu1_8+\mu1_4^2,\\
W_1 &= 2\mu1_1\mu1_{17}-\mu1_{11}^2-2\mu1_{11}\mu1_{15}+\mu1_{13}^2+2\mu1_{13}\mu1_5-\mu1_{15}^2-2\mu1_{17}\mu1_9+\mu1_5^2.
\end{align}

\begin{align}
V_0 &= \mu1_0^2-\mu1_8^2-2\mu1_0\mu1_{16}+\mu1_{10}^2+2\mu1_{10}\mu1_{14}-\mu1_{12}^2-2\mu1_{12}\mu1_4+\mu1_{14}^2+2\mu1_{16}\mu1_8-\mu1_4^2,\\
V_1 &= \mu1_1^2-\mu1_9^2-2\mu1_1\mu1_{17}+\mu1_{11}^2+2\mu1_{11}\mu1_{15}-\mu1_{13}^2-2\mu1_{13}\mu1_5+\mu1_{15}^2+2\mu1_{17}\mu1_9-\mu1_5^2.
\end{align}

\section*{Nota su notazione $d_J,d_K$}
Se usi la convenzione con \(d_J,d_K\), in questo output valgono le identificazioni:
\begin{equation}
d_J \equiv d_1, \qquad d_K \equiv d_2.
\end{equation}

\end{document}
